\section{问题二的模型与求解}

为选择第二个检测点，应先对每个候选位置分析可能出现的反馈，确定各反馈下的剩余可行区域，并根据能否原地完成光学精确定位计算相应损失。在此基础上，我们建立了两种选点模型。\textbf{基于贝叶斯更新的期望定位损失模型}先根据首次观测更新未知状态的概率，再预测各反馈并对相应损失加权平均，选择平均损失最小的检测点，侧重提高所设概率假设下的平均定位效果；\textbf{最小化最坏剩余半径模型}则考察每个候选点可能遭遇的最不利反馈，选择最坏损失最小的检测点，侧重控制不利情形下仍未消除的位置不确定性。两种模型采用相同的区域构造与损失规则，区别在于选点时依据平均表现，还是依据最坏情形下的保证。

\subsection{准备工作：由观测信息构造定位区域}

每次反馈先转化为源位置约束，再与已有信息取交。利用旋转对称性，取 \(s_1=(a,0)\)、\(\theta_1=\beta\)，其中 \(a\ge0\)、\(\beta\in[-\pi,\pi)\) 为给定参数。考察第二点 \(q\) 及其可能读数 \(\theta\) 时，令 \(s_2=q\)、\(\theta_2=\theta\)，并简记 \(d_1(g)=\|g-s_1\|\)、\(d(g;q)=\|g-q\|\)。

根据题意，已知某次观测获得了正常示向度，这就同时给出方向、有效接收和非强信号三条信息，综合可知源位置可行区域为 \(K_1=\{g\in\Omega\cap W_1:r_{\mathrm{strong}}<d_1(g)\le1500\}\)。对任一 \(g\in K_1\)，固定的接收半径满足 \(\max\{1000,d_1(g)\}\le R\le1500\)。

在候选点 \(q\) 进行第二次检测，可能得到强信号、无信号或正常示向度，分别记为 \(\mathrm H\)、\(\mathrm N\) 和 \(\theta\)。每种反馈都能进一步排除首次可行区域 \(K_1\) 中的部分位置。

若收到强信号，说明源距第二点不超过5米；若没有收到信号，由于首次能够接收且接收半径固定，说明第二点比首次检测点离源更远，并且源距超过接收半径的下限1000米；若得到正常示向度，则源位于新读数对应的角域内，且源距大于5米、不超过1500米。因此，三种反馈对应的剩余可行区域分别为
\begin{equation}
\begin{aligned}
K_{\mathrm H}(q)&=\{g\in K_1:d(g;q)\le r_{\mathrm{strong}}\},\\
K_{\mathrm N}(q)&=\{g\in K_1:d(g;q)>\max\{1000,d_1(g)\}\},\\
K_\theta(q)&=\{g\in K_1\cap W_2:r_{\mathrm{strong}}<d(g;q)\le1500\}.
\end{aligned}
\label{eq:q2-pair-regions}
\end{equation}
其中，\(W_2\) 是以第二点 \(q\) 为顶点、以示向度 \(\theta\) 为中心方向的测向角域。若某种反馈对应的区域为空，说明它与已有信息不相容，不可能出现。将其余可能出现的反馈组成的集合记为 \(F(q)\)，供后续评价候选点时使用。

为衡量剩余的位置不确定性，用能够覆盖整个可行区域 \(K\) 的最小圆的半径表示区域大小，记为 \(\rho(K)\)；用 \(d_{\max}(q;K)\) 表示检测点 \(q\) 到该区域的最远距离。第二次反馈为 \(z\) 时，剩余区域为 \(K_z(q)\)，其最小包围圆半径简记为 \(\rho_z(q)\)。前者用于评价定位精度，后者用于判断能否在当前点完成光学定位。

第二个检测点不宜离首次可行区域太远。若两者相距超过1500米，无论干扰源位于区域内何处，都只能得到无信号反馈，无法进一步缩小定位范围。因此，只考虑距首次可行区域不超过1500米、且不同于首次检测点的位置，将这一候选范围记为 \(Q\)。检测点可以位于目标圆域 \(\Omega\) 之外。

\subsection{反馈后的统一定位损失}

评价时点取为第二次无线电检测后、完成能够在当前点保证成功的光学精确定位之时。若整个可行区域都在 \(q\) 的光学范围内，则可确定源的准确位置，损失为零；否则以剩余区域的最小包围圆半径衡量位置不确定性。对每个可实现反馈，定义以米计的损失
\begin{equation}
 L_z(q)=
 \begin{cases}
 0,&d_{\max}(q;K_z(q))\le r_{\mathrm{opt}},\\
 \rho_z(q),&d_{\max}(q;K_z(q))>r_{\mathrm{opt}}.
 \end{cases}
\label{eq:q2-pair-loss}
\end{equation}
强信号必能原地定位，故 \(L_{\mathrm H}=0\)；无信号对应的位置距 \(q\) 超过1000米，故 \(L_{\mathrm N}=\rho_{\mathrm N}\)。正常示向度也可能使整个剩余区域落入当前点的20米范围，此时同样取零损失。这里检验的是区域到当前点的最远距离，仅有 \(\rho_z\le20\) 只能保证存在一个可定位位置，不能保证在 \(q\) 原地成功。无信号虽未提供新角度，其距离约束仍可能缩小区域，该收益也通过 \(\rho_{\mathrm N}\) 计入。

\subsection{基于贝叶斯更新的期望定位损失模型}

本模型采用贝叶斯决策方法，按照“先验描述—观测更新—反馈预测—损失最小化”的顺序选择第二检测点。先验描述观测前对未知状态的认识，首次观测用于调整各个可能状态的权重，更新后的分布则用于评价下一次检测的预期收益。

\subsubsection{固定未知量的贝叶斯描述}

源位置、接收半径和各检测地点的误差在所考虑的任务时段内均有固定取值，但在观测前并不为我们所知。\textbf{贝叶斯模型中的概率分布描述对这些固定未知量的认识。} 随着观测信息增加，我们对其可能取值的判断发生变化，各未知量本身仍保持固定。因此，对地点误差设置概率分布，与题目规定的“同一地点重复检测不改变误差”并不矛盾。

按照前文的模型假设，对源位置 \(G\) 采用 \(\Omega\) 内按面积均匀的先验，对固定接收半径采用 \(R\sim\operatorname{U}[1000,1500]\)，对两个不同检测地点的误差采用 \(E_1,E_2\sim\operatorname{U}[-\alpha,\alpha]\)，并在观测前假设 \(G,R,E_1,E_2\) 相互独立。均匀性与独立性是计算反馈概率所采用的额外假设，不能由题目给出的取值范围直接推出；尤其是不同地点的误差，也可能存在空间相关性。

同一地点的多次检测共享同一个误差 \(E_i\)，不能作为新的独立样本，也不能通过取平均缩小误差。本模型只在 \(q\ne s_1\) 的新地点预测第二次反馈，并在更新后保留未知量之间产生的关联。所得平均定位效果以所设先验为前提，后文另用不依赖概率分布的最坏情形模型给出定位保证。

\subsubsection{首次观测后的概率更新}

将首次正常示向度信息记为 \(I_1=(s_1,\beta)\)。贝叶斯更新将观测前的先验权重与观测似然相乘，再归一化得到后验；其中，似然反映各候选状态产生当前观测的可能程度。

对于符合方向约束的位置，均匀角误差给出的似然密度相同，但正常接收还要求 \(R\ge d_1(g)\)。因此，需要先计算距离为 \(d\) 时的接收概率
\begin{equation}
 w(d)=\Pr(R\ge d)=
 \begin{cases}
 1,&d\le1000,\\
 (1500-d)/500,&1000<d<1500,\\
 0,&d\ge1500.
 \end{cases}
\label{eq:q2-pair-weight}
\end{equation}
据此将候选位置的权重归一化，得到首次位置后验
\begin{equation}
 p_1(g)=\frac{\mathbf1_{K_1}(g)w(d_1(g))}{C_1},
 \qquad C_1=\int_{K_1}w(d_1(g))\,\mathrm dg>0.
\label{eq:q2-pair-posterior}
\end{equation}
这里 \(C_1\) 保证密度积分为1。该更新的直观含义是，较远的位置需要较大的接收半径才能解释首次接收，因而在更新后获得较低权重。进一步，对于正后验密度的位置，有 \(R\mid G=g,I_1\sim\operatorname{U}[\max\{1000,d_1(g)\},1500]\)；位置与半径在后验中由此产生关联，第二次预测须保留这一关系。

\subsubsection{第二次反馈的概率预测}

在首次观测的后验基础上预测第二次反馈，将相应的反馈条件加入首次观测条件，并对未知半径积分。记 \(d_\vee(g;q)=\max\{d_1(g),d(g;q)\}\)。两次正常接收要求 \(R\ge d_\vee\)，对应权重为 \(w(d_\vee)\)；先接收而后无信号的权重则为 \(w(d_1)-w(d_\vee)\)。结合方向信息，三类反馈的未归一化位置权重为
\begin{equation}
\begin{aligned}
 b_{\mathrm H}(g;q)&=\mathbf1_{K_{\mathrm H}(q)}(g)w(d_1(g)),\\
 b_{\mathrm N}(g;q)&=\mathbf1_{K_1}(g)\bigl[w(d_1(g))-w(d_\vee(g;q))\bigr],\\
 b_\theta(g;q)&=\frac{\mathbf1_{K_\theta(q)}(g)}{2\alpha}w(d_\vee(g;q)).
\end{aligned}
\label{eq:q2-pair-feedback-weights}
\end{equation}
在不同检测点误差先验独立的假设下，正常分支的 \(1/(2\alpha)\) 为新地点误差的条件预测密度。对上述权重积分并归一化，即得三类反馈的预测概率或密度
\begin{equation}
\begin{aligned}
 P_{\mathrm H}(q)&=\frac{1}{C_1}\int_{K_1}b_{\mathrm H}(g;q)\,\mathrm dg,
 &P_{\mathrm N}(q)&=\frac{1}{C_1}\int_{K_1}b_{\mathrm N}(g;q)\,\mathrm dg,\\
 f_\Theta(\theta;q)&=\frac{1}{C_1}\int_{K_1}b_\theta(g;q)\,\mathrm dg.
\end{aligned}
\label{eq:q2-pair-prediction}
\end{equation}
前两项为离散反馈概率，最后一项为正常角度的预测密度分量，其在 \([0,2\pi)\) 上的积分等于正常反馈发生的概率，三者合计为1。

\subsubsection{期望损失与贝叶斯决策}

以 \(Z\) 表示尚未出现的第二次反馈，按式~\eqref{eq:q2-pair-prediction} 的预测概率对统一损失求平均，得到
\begin{equation}
 J_{\mathrm E}(q)=\mathbb E[L_Z(q)\mid I_1,q]
 =P_{\mathrm N}(q)\rho_{\mathrm N}(q)
 +\int_0^{2\pi}f_\Theta(\theta;q)L_\theta(q)\,\mathrm d\theta.
\label{eq:q2-pair-expected-objective}
\end{equation}
强信号以零损失参与加权，正常分支的光学定位收益已包含在 \(L_\theta\) 中。概率或预测密度为零的反馈贡献按零计，无须定义空区域的包围圆。最小化 \(J_{\mathrm E}(q)\) 即得到当前先验与损失定义下的贝叶斯选点决策。该指标衡量平均剩余不确定性，其较小并不意味着每种反馈下的损失都较小。

\subsection{最小化最坏剩余半径的模型}

为控制不利反馈下的剩余不确定性，对同一损失取全部可实现反馈中的上确界，得到
\begin{equation}
 J_{\mathrm W}(q)=\sup_{z\in F(q)}L_z(q).
\label{eq:q2-pair-worst-objective}
\end{equation}
该指标不依赖前述概率假设，保留全部符合题意的反馈，包括边界状态及退化交集；即使某些状态在连续先验下概率为零，也不能从最坏情形中排除。在期望模型适用时，对同一点有 \(J_{\mathrm E}(q)\le J_{\mathrm W}(q)\)，但分别最小化两者可能得到不同选点，体现平均效果与最坏保证之间的取舍。

\subsection{选点与候选区域}

给定首次观测参数 \((a,\beta)\)，对任一准则 \(J\in\{J_{\mathrm E},J_{\mathrm W}\}\)，求 \(J^*=\inf_{q\in Q}J(q)\)，最小值可达时取 \(q^*\in\arg\min_{q\in Q}J(q)\)。为输出一片较好的候选区域，给定容差 \(\epsilon>0\)，保留 \(Q_{J,\epsilon}=\{q\in Q:J(q)\le J^*+\epsilon\}\)。其中 \(\epsilon\) 以米计，表示相对于相应最优指标允许增加的损失，最优值为零时仍适用。

两种模型均可采用检测点空间搜索与局部加密：对每个候选点构造反馈区域、计算包围圆并检查原地定位条件，期望准则对正常角度积分，最坏准则搜索其损失上确界，再分别合并离散反馈的贡献。数值所得候选区域须结合积分或角度搜索误差及外层搜索精度核查。
